% Part: proof-theory
% Chapter: sequent-calculus
% Section: rules-G3c

\documentclass[../../../include/open-logic-section]{subfiles}

\begin{document}

\begin{table}
\begin{tabular}{@{}c@{}c@{}}
\multicolumn{2}{@{}c@{}}{अभिगृहीतियाँ}\\[2ex]
$!C, \Gamma \Sequent !C$ & $\lfalse, \Gamma \Sequent \Delta$
\\[2ex]
\multicolumn{2}{@{}c@{}}{तार्किक नियम}\\[2ex]
\Axiom$\lnot !A, \Gamma \fCenter !A $
\RightLabel{\LeftR{\lnot}}
\UnaryInf$\lnot !A, \Gamma \fCenter \Delta$
\DisplayProof
&
\Axiom$!A, \Gamma \fCenter$
\RightLabel{\RightR{\lnot}}
\UnaryInf$ \Gamma \fCenter \lnot !A $
\DisplayProof
\\[4ex]
\Axiom$ !A, !B, \Gamma \fCenter \Delta$
\RightLabel{\LeftR{\land}}
\UnaryInf$ !A \land !B, \Gamma \fCenter \Delta$
\DisplayProof
&
\Axiom$\Gamma \fCenter !A$
\Axiom$ \Gamma \fCenter !B$
\RightLabel{\RightR{\land}}
\BinaryInf$ \Gamma \fCenter !A \land !B $
\DisplayProof
\\[4ex]\Axiom$!A, \Gamma \fCenter \Delta$
\Axiom$!B, \Gamma \fCenter \Delta$
\RightLabel{\LeftR{\lor}}
\BinaryInf$!A \lor !B, \Gamma \fCenter \Delta$
\DisplayProof
&
\Axiom$\Gamma \fCenter !A, !B$
\RightLabel{\RightR{\lor}}
\UnaryInf$ \Gamma \fCenter !A \lor !B$
\DisplayProof
\\[4ex]
\Axiom$!A \lif !B, \Gamma \fCenter !A$
\Axiom$!B, \Gamma \fCenter \Delta$
\RightLabel{\LeftR{\lif}}
\BinaryInf$!A \lif !B, \Gamma \fCenter \Delta$
\DisplayProof
&
\Axiom$!A, \Gamma \fCenter !B$
\RightLabel{\RightR{\lif}}
\UnaryInf$ \Gamma \fCenter !A \lif !B $
\DisplayProof
\\[4ex]
\Axiom$ !A(t), \lforall[x][!A(x)]\Gamma \fCenter \Delta$
\RightLabel{\LeftR{\lforall}}
\UnaryInf$ \lforall[x][!A(x)],\Gamma \fCenter \Delta$
\DisplayProof
&
\Axiom$ \Gamma \fCenter !A(a) $
\RightLabel{\RightR{\lforall}}
\UnaryInf$ \Gamma \fCenter \lforall[x][!A(x)]$
\DisplayProof
\\[4ex]
\Axiom$ !A(a), \Gamma \fCenter \Delta $
\RightLabel{\LeftR{\lexists}}
\UnaryInf$ \lexists[x][!A(x)], \Gamma \fCenter \Delta$
\DisplayProof
&
\Axiom$ \Gamma \fCenter !A(t) $
\RightLabel{\RightR{\lexists}}
\UnaryInf$ \Gamma \fCenter \lexists[x][!A(x)]$
\DisplayProof
\end{tabular}
\caption{अंतःप्रज्ञात्मक तर्क के लिए \Log{G3i} के नियम।
$\RightR{\forall}$ और $\LeftR{\exists}$ में !!{constant}~$a$ निष्कर्ष
अनुक्रम में नहीं आना चाहिए। $\Gamma$, !!{formula}s का बहुसमुच्चय है और
$\Delta$ या तो रिक्त हो सकता है या उसमें एक अकेला !!{formula} हो सकता है।
अभिगृहीतियों में !!{formula}~$!C$ परमाण्विक होना चाहिए। \Log{G1m},
\Log{G1i} में से $\lfalse, \Gamma \Sequent \Delta$ अभिगृहीतियाँ हटाने पर
मिलता है।}
\ollabel{tab:G3c}
\end{table}

\end{document}
